\par\bigskip
\begin{center}\pdfbookmark[2]{7. Säkulare und periodische Störungen}{ceres-007}[7.]\end{center}
\noindent\hspace*{\parindent}Aus obigen Formeln folgt nunmehr

I) Säkularer Zuwachs der Exzentrizität \(=\tfrac12e'a\mathfrak C^{(0)}nt\sin(\tilde\omega'-\tilde\omega)\). [Nun ist] \(\mathfrak C^{(0)}=\tfrac12Q''\), also der Zuwachs \(=\tfrac14e'aQ''nt\sin(\tilde\omega'-\tilde\omega)\)
\[
\begin{array}{cc}
\text{Für Ceres ist:}&\text{[Für Jupiter:]}\\[.5ex]
\tilde\omega=352^{0}\,58'\,0''&\tilde\omega'=191^{0}\,10'\,20''\\[.5ex]
\log n\ \text{[täglich]}=2{,}886\,2959&\log e'=8{,}681\,3317\,[-10]
\end{array}
\]
[und hieraus findet man
\[
\text{jährliche Abnahme der Exzentrizität]}\quad0{,}00000\,59089.
\]

II) Säkulare Zunahme der Länge des Aphels:
\[
=\frac a4\left(B^{(0)}-2\mathfrak B^{(0)}-2\frac{e'}e\mathfrak C^{(0)}\cos(\tilde\omega'-\tilde\omega)\right)nt.
\]
Nun ist
\[
B^{(0)}-2\mathfrak B^{(0)}=Q',
\]
also jene Zunahme
\[
=\frac a4\left(Q'-\frac{e'}eQ''\cos(\tilde\omega'-\tilde\omega)\right)nt.
\]

Für Ceres, jährliche Bewegung der Apsiden \(=+70''{,}1503.\)

Nach diesen Säkularänderungen der Exzentrizität und des Aphels fallen die oben rot\refstepcounter{footnote}\hyperlink{ceres7-note}{[*)]}\begingroup
\makeatletter
\long\def\@makefntext#1{\noindent[*) #1}
\makeatother
\footnotetext{\hypertarget{ceres7-note}{}Seite 389, wo sie auch im Abdruck unterstrichen sind.]}
\endgroup{} unterstrichenen Gleichungen ganz weg.

III) Der konstante Teil von \(\delta r\) ist \(=[-aa(2g+\mathfrak A^{(0)})=]+\tfrac16aaB^{(0)}\)
\[
\text{[also]}\quad=-0{,}000\,9475.
\]

IV) Der Koeffizient von \(\cos iD\) [in \(\delta r\)] wird (weil \(i\) oben sowohl für positive als auch für negative Zahlen stand und nun auf die positiven beschränkt wird)
\[
\left[\frac{2aa\mathfrak A^{(i)}}{hh-1}=\right]\frac{2aa}{hh-1}\left(\frac ihA^{(i)}+\tfrac12B^{(i)}\right)\mathpunct{\cdot}
\]

\hypertarget{native-398}{}[In unserem Fall, \(\log n'=2{,}475\,8562\),] ergibt sich \(\log\theta=9{,}589\,5603\) [und damit der Teil von \(\delta r\):]
\[
\begin{aligned}
&+0{,}001\,0304\cos D&&-0{,}000\,0350\cos5D\\
&-0{,}003\,8023\cos2D&&-0{,}000\,0128\cos6D\\
&-0{,}000\,4206\cos3D&&-0{,}000\,0050\cos7D.\\
&-0{,}000\,1077\cos4D
\end{aligned}
\]

V) In \(\delta v\) wird der Koeffizient von \(\sin iD\)
\[
\begin{aligned}
&-\frac{4ai(1-\theta)}{ii(1-\theta)^2-1}\left(\frac1{1-\theta}A^{(i)}+\tfrac12B^{(i)}\right)\\
&+\frac a{i(1-\theta)}\left(\frac{3A^{(i)}}{1-\theta}+2B^{(i)}\right)
\end{aligned}
\]
[also in unserem Fall die entsprechenden Glieder:]
\[
\begin{aligned}
&-231''{,}94\sin D&&+3''{,}05\sin5D\\
&+496{,}68\sin2D&&+1{,}07\sin6D\\
&+44{,}15\sin3D&&+0{,}41\sin7D.\\
&+10{,}07\sin4D
\end{aligned}
\]

VI) Zur Bestimmung der von \(M\) und \(M'+D\) [abhängigen Glieder] dient folgende Untersuchung. [Bisher sind in \(\delta v\) die Integrationskonstanten aus Gleichung 11 unberücksichtigt geblieben. Sie können beliebig gewählt werden, da sie sich mit den elliptischen Elementen \(e\) und \(\tilde\omega\) vereinigen, und sollen so bestimmt werden, dass die von \(M\) und \(D+M'=M+\tilde\omega-\bar{\tilde\omega}'\) abhängigen Glieder in \(\delta v\) verschwinden. Werden die Integrationskonstanten \(f_1\) und \(f_1'\) so eingeführt, dass die Konstante von \(\int rQ\sin V\mathord{.}\,n\,\mathrm{d}t\) gleich \(aaef_1+aae'f_1'\cos(\tilde\omega-\tilde\omega')\) und die aus \(\int rQ\cos V\mathord{.}\,n\,\mathrm{d}t\) entstehende Konstante gleich \(aae'f_1'\sin(\tilde\omega-\tilde\omega')\) wird, so ist der entsprechende Teil von \(\delta r\) unter Vernachlässigung der zweiten Potenzen der Exzentrizitäten
\[
\alpha)\qquad aaef_1\cos M+aae'f_1'\cos(M+\tilde\omega-\tilde\omega')
\]
und dadurch entsteht in \(\delta v\) der Teil
\[
\beta)\qquad-2aef_1\sin M-2ae'f_1'\sin(M+\tilde\omega-\tilde\omega').]
\]

Die [gesamten von \(M\) und \(D+M'\) abhängigen] Glieder einschließlich derer [\(\alpha\)], die sich aus den Konstanten ergeben, seien
\hypertarget{native-399}{}
\[
\text{in }\delta r\ .\, .\, .\, .\quad aaef\cos M+aae'f'\cos(M+\tilde\omega-\tilde\omega')
\]
[sodass also
\[
\begin{aligned}
f&=f_1-\tfrac1{24}B^{(0)}-\tfrac14\mathfrak B^{(0)}\\
f'&=f_1'-\tfrac14\mathfrak C^{(0)}.
\end{aligned}
\]

Damit ergibt sich der gesamte von den beiden Argumenten abhängige Teil
\[
\begin{aligned}
\text{in }\delta v\ .\, .\, .\, .\quad&-2a\{f-\tfrac16B^{(0)}-\tfrac12\mathfrak B^{(0)}+\tfrac14\beta^{(0)}\}e\sin M\\
&-2a\{f'-\tfrac12\mathfrak C^{(0)}+\tfrac14\gamma^{(0)}\}e'\sin(M+\tilde\omega-\tilde\omega').]
\end{aligned}
\]
Man bestimmt \(f,f'\) so, dass die Summe dieser Größen in \(\delta v\) verschwindet. So bleibt
\[
\begin{aligned}
\text{in }\delta r\ .\, .\, .\, .\quad&aae(\tfrac12\mathfrak B^{(0)}-\tfrac14\beta^{(0)}+\tfrac16B^{(0)})\cos M\\
&+aae'(\tfrac12\mathfrak C^{(0)}-\tfrac14\gamma^{(0)})\cos(M+\tilde\omega-\tilde\omega')
\end{aligned}
\]
oder
\[
\begin{aligned}
&aae(\tfrac5{12}B^{(0)}+\tfrac14C^{(0)})\cos M\\
&+aae'(\tfrac14A'-\tfrac14C')\cos(M+\tilde\omega-\tilde\omega')
\end{aligned}
\]
[und numerisch]
\[
-0{,}000\,0526\cos M\qquad+0{,}000\,0141\cos(M+\tilde\omega-\tilde\omega').
\]

[Diese beiden Glieder lassen sich zu einem zusammenfassen. Bezeichnet man sie vorübergehend mit
\[
b_1\cos M+b_2\cos(M+\tilde\omega-\tilde\omega')
\]
und bestimmt man die Konstanten \(c\) und \(C\) aus den Gleichungen
\[
\begin{aligned}
c\cos C&=b_1+b_2\cos(\tilde\omega-\tilde\omega')\\
c\sin C&=b_2\sin(\tilde\omega-\tilde\omega'),
\end{aligned}
\]
so erhalten sie die Form
\[
c\cos(M+C),
\]
numerisch] \(-0{,}000\,0633\cos(M+25^{0}\,56'\,46'').\)

VII) In \(\delta r\) wird der Koeffizient von \(\cos(iD+M)\), [der mit] \(aaeE^{(i)}\) [bezeichnet werden mag]
\[
\begin{aligned}
aaeE^{(i)}&=aae\left\{\frac{\mathfrak B^{(i)}}{hh+2h}+\tfrac12\mathfrak A^{(i)}\left(\frac1{h+2}-\frac{4h}{hh-1}+\frac3h\right)\right\}\\
&=\frac{aae}{h(h+2)}\left\{\begin{aligned}
&\tfrac12C^{(i)}-\tfrac12B^{(i)}-\frac{3A^{(i)}}{1-\theta}\\
&-2\mathfrak A^{(i)}\left(\frac{3+ih(h-1)}{hh-1}\right)
\end{aligned}\right\}\mathpunct{\cdot}
\end{aligned}
\]

VIII) In \(\delta r\) wird der Koeffizient von \(\cos((i+1)D+M')\), [der mit] \(aae'D^{(i)}\) \hypertarget{native-400}{}[bezeichnet werden mag]
\[
\begin{aligned}
aae'D^{(i)}&=\frac{aae'\mathfrak C^{(i)}}{hh+2h}\\
&=\frac{aae'}{hh+2h}\left\{\begin{aligned}
&\frac{i+1}{h+1}\bigl((2i+1)A^{(i+1)}-B^{(i+1)}\bigr)\\
&+(i+\tfrac12)B^{(i+1)}-\tfrac12C^{(i+1)}
\end{aligned}\right\}\mathpunct{\cdot}
\end{aligned}
\]

[Auch hier unterscheiden sich die Argumente \(iD+M\) und \((i+1)D+M'\) nur um konstante Größen. Die Glieder lassen sich ebenfalls zu einem zusammenfassen. Setzt man ähnlich wie unter VI]
\[
\alpha\genfrac{}{}{0pt}{1}{\sin}{\cos}(x-A)+\beta\genfrac{}{}{0pt}{1}{\sin}{\cos}(x-B)
=f\genfrac{}{}{0pt}{1}{\sin}{\cos}(x-A+z),
\]
[so bestimmen sich \(f\) und \(z\) aus
\[
\begin{aligned}
\beta\sin(A-B)&=f\sin z\\
\alpha+\beta\cos(A-B)&=f\cos z.]
\end{aligned}
\]

IX) In \(\delta v\) wird der Koeffizient von \(\sin(iD+M)\)
\[
\begin{aligned}
&=ae\left\{\begin{aligned}
&-2(h+1)E^{(i)}+\frac{2\mathfrak B^{(i)}-\tfrac12\beta^{(i)}}{h+1}\\
&-(2h+1)\frac{\mathfrak A^{(i)}}{hh-1}
\end{aligned}\right\}\\
&=\frac{ae}{h+1}\left\{\begin{aligned}
&-2E^{(i)}-2\mathfrak A^{(i)}\frac{\tfrac12i(h-1)-3}{hh-1}\\
&+\frac{(i+1)A^{(i)}}{1-\theta}
\end{aligned}\right\}\mathpunct{\cdot}
\end{aligned}
\]

X) In \(\delta v\) wird der Koeffizient von \(\sin((i+1)D+M')\)
\[
\begin{aligned}
&=ae'\left\{-2(h+1)D^{(i)}+\frac{2\mathfrak C^{(i)}-\tfrac12\gamma^{(i)}}{h+1}\right\}\\
&=\frac{ae'}{h+1}\left\{\begin{aligned}
&-2D^{(i)}\\
&-\tfrac12\frac{i+1}{h+1}\bigl((2i+1)A^{(i+1)}-B^{(i+1)}\bigr)
\end{aligned}\right\}\mathpunct{\cdot}
\end{aligned}
\]
\par\medskip\centerline{\rule{.23\textwidth}{.35pt}}\medskip
